\subsection{R 中可和的正有限数族}

定理 1. 族 \((x_{t})\) （由有限实数 \(\geqslant 0\) 组成）在 \(\mathbf{R}\) 中可和当且仅当该族的有限部分和组成的集合在 \(\mathbf{R}\) 中上有界。此时，该集合的上确界就是族 \((x_{t})\) 的和。

对指标集 I 的每个有限子集 H，令 \(s_{H} = \sum_{i \in H} x_{i}\) ；由于诸 \(x_{i}\) \(\geqslant 0\) ，关系 \(H \subset H'\) 蕴含 \(s_{H} \leqslant s_{H'}\) 。换句话说，映射 \(H \to s_{H}\) 是 I 的有限子集组成的有向集 \(\mathfrak{F}(I)\) 上的递增函数；因此（§ 5，第 2 节，定理 2 的推论）它有有限极限当且仅当它上有界。

注. 设 \((\mathrm{H}_{\lambda})\) 是 I 的有限子集的一个族，使得对 I 的每个有限子集 H，存在指标 \(\lambda\) 满足 \(H \subset H_{\lambda}\) ；则 \((x_{t})\) 可和当且仅当诸 \(s_{H_{1}}\) 组成的族在 R 中上有界。特别地，设 \((x_{n})\) 是有限实数 \(\geqslant 0\) 的序列，且对每个整数 n 令 \(s_{n} = \sum_{p=0}^{n} x_{p}\) ；则序列 \((x_{n})\) 在 R 中可和当且仅当对一个严格递增整数序列 \((n_{k})\) ，部分序列 \((s_{n_{k}})\) 在 R 中上有界。

例. 1) 对每个数 \(q\) （满足 \(0 \leqslant q < 1\) ），序列 \((q^n)\) （'公比为 \(q''\) 的等比数列'）在 \(\mathbf{R}\) 中可和，因为 \(s_n = \frac{1 - q^{n+1}}{1 - q} \leqslant \frac{1}{1 - q}\) ；该序列的和为 \(\lim_{n \to \infty} s_n = \frac{1}{1 - q}\) 。

\begin{enumerate}
\def\labelenumi{\arabic{enumi})}
\setcounter{enumi}{1}
\tightlist
\item
  设 \(a\) 与 \(b\) 是两个数，满足 \(0 \leqslant a < 1\) 与 \(0 \leqslant b < 1\) ；则族 \((a^{m}b^{n})_{(m,n) \in \mathbb{N} \times \mathbb{N}}\) 在 \(\mathbb{R}\) 中可和。因为 \(\mathbf{N} \times \mathbf{N}\) 的每个有限子集都含于形如 \([0, p] \times [0, p]\) 的子集，且我们有
\end{enumerate}

\[
\sum_ {m = 0} ^ {p} \sum_ {n = 0} ^ {p} a ^ {m} b ^ {n} = \left(\sum_ {m = 0} ^ {p} a ^ {m}\right) \left(\sum_ {n = 0} ^ {p} b ^ {n}\right) = \frac {1 - a ^ {p + 1}}{1 - a}. \frac {1 - b ^ {p + 1}}{1 - b} \leqslant \frac {1}{(1 - a) (1 - b)}.
\]

\begin{enumerate}
\def\labelenumi{\arabic{enumi})}
\setcounter{enumi}{2}
\tightlist
\item
  对每个整数 \(p > 1\) ，序列 \((n^{-p})\) ( \(n > 0\) )可和，因为
\end{enumerate}

\[
s _ {2 ^ {n + 1}} - s _ {2 ^ {n}} = \sum_ {k = 1} ^ {2 ^ {n}} (2 ^ {n} + k) ^ {- p} <   2 ^ {n}. (2 ^ {n}) ^ {- p}
\]

从而把这些不等式相加，得

\[
s _ {2 ^ {n}} <   \frac {1}{1 - 2 ^ {1 - p}}.
\]

\begin{enumerate}
\def\labelenumi{\arabic{enumi})}
\setcounter{enumi}{3}
\tightlist
\item
序列 \((1 / n)\) \((n > 0)\) 在 \(\mathbf{R}\) 中不可和，因为
\end{enumerate}

\[
s _ {2 ^ {n + 1}} - s _ {2 ^ {n}} = \sum_ {k = 1} ^ {2 ^ {n}} \frac {1}{2 ^ {n} + k} > \frac {2 ^ {n}}{2 ^ {n + 1}} = \frac {1}{2}
\]

从而把这些不等式相加，得

\[
s _ {2 ^ {n}} > n / 2
\]

故定理 1 的判别准则不满足。

\begin{enumerate}
\def\labelenumi{\arabic{enumi})}
\setcounter{enumi}{4}
\tightlist
\item
  设 \((\mathbf{I}_n)\) 是两两不相交的非空开区间的一个序列，它们都含于一个长度为有限数 \(l\) 的区间中。该族中有限个区间的长度之和 \(\leqslant l\) （§ 1，第 5 节），因此诸 \(\mathbf{I}_n\) 的长度组成的族在 \(\mathbf{R}\) 中可和，且其和 \(\leqslant l\) 。
\end{enumerate}

定理 2（比较原理）. 设 \((x_{\iota})_{\iota\in I}\) 与 \((y_{\iota})_{\iota\in I}\) 是有限实数 \(\geqslant 0\) 的两个族，满足 \(x_{\iota}\leqslant y_{\iota}\) 对一切 \(\iota\) 成立。若 \((y_{\iota})\) 在 \(\mathbb{R}\) 中可和，则 \((x_{\iota})\) 也可和，且我们有 \(\sum_{\iota}x_{\iota}\leqslant\sum_{\iota}y_{\iota}\) ；此外若存在指标 \(\chi\) 使 \(x_{\chi}<y_{\chi}\) ，则 \(\sum_{\iota}x_{\iota}<\sum_{\iota}y_{\iota}\) 。

假设蕴含

\[
\sum_ {\iota \in \mathbf {H}} x _ {\iota} \leqslant \sum_ {\iota \in \mathbf {H}} y _ {\iota},
\]

对 I 的每个有限子集 H 成立，定理的第一部分由此得证；关于和的不等式则由不等式延拓原理（§ 5，第 2 节，定理 1）得出。若 \(x_{\chi} < y_{\chi}\) ，则

\[
\sum_ {\iota} x _ {\iota} = x _ {\chi} + \sum_ {\iota \neq \chi} x _ {\iota} <   y _ {\chi} + \sum_ {i \neq x} y _ {\iota} = \sum_ {\iota} y _ {\iota}.
\]

这个定理提供了判断序列 \((x_{n})\) （由实数 \(\geqslant0\) 组成）在 R 中是否可和的最常用判别准则：我们设法把给定的序列与一个较简单的序列 \((y_{n})\) 相比较，后者的可和性我们已经知道。如果存在有限数 a \textgreater{} 0 ，使得从某个点起对一切 n 有 \(x_{n} \leqslant ay_{n}\) ，且 \((y_{n})\) 可和，则 \((x_{n})\) 可和；反之，如果存在有限数 b \textgreater{} 0 ，使得从某个点起对一切 n 有 \(x_{n} \geqslant by_{n}\) ，且 \((y_{n})\) 在 R 中不可和，则 \((x_{n})\) 在 R 中不可和。我们将在后续分册中看到，在最常出现的情形如何得到这样的比较序列。

例. 1) 设 \(a\) 是有限实数 \(>0\) ，考虑序列 \(\left(\frac{a^n}{n!}\right)\) ；设 \(n_0\) 是使 \(a < n_0\) 的最小整数。则

对每个 \(n \geqslant n_0\) 我们有

\[
\frac {a ^ {n}}{n !} \leqslant \frac {a ^ {n _ {0}}}{n _ {0} !} \cdot \left(\frac {a}{n _ {0}}\right) ^ {n - n _ {0}};
\]

由于 \(q = \frac{a}{n_0} < 1\) ，序列 \((q^{n - n_0})\) 可和，因此序列 \(\left(\frac{a^n}{n!}\right)\) 也可和。

\begin{enumerate}
\def\labelenumi{\arabic{enumi})}
\setcounter{enumi}{1}
\tightlist
\item
  设 \((a_{n})\) 是正数的一个可和序列。由于
\end{enumerate}

\[
\lim _ {n \rightarrow \infty} a _ {n} = 0,
\]

存在整数 \(n_0\) ，使得 \(a_n \leqslant 1\) 只要 \(n \geqslant n_0\) 。于是对每个 \(n \geqslant n_0\) 有 \(a_n^2 \leqslant a_n\) ，因此序列 \((a_n^2)\) 在 \(\mathbb{R}\) 中可和。序列 \((a_n^p)\) 对每个整数 \(p > 1\) 也可和。

\begin{enumerate}
\def\labelenumi{\arabic{enumi})}
\setcounter{enumi}{2}
\tightlist
\item
  设 \(a\) 与 \(b\) 是两个数 \(< 1\) ；则
\end{enumerate}

\[
\frac {1}{a ^ {m} + b ^ {n}} \leqslant \frac {1}{2 (\sqrt {a}) ^ {m} (\sqrt {b}) ^ {n}}
\]

因而族 \(\left(\frac{1}{a^m + b^n}\right)\) 在 \(\mathbf{R}\) 中可和。

推论. 设 \((x_{i})_{i\in I}\) 是 R 中有限数 \(\geqslant 0\) 的可和族。若 H 是 I 的任一子集，我们有

\[
\sum_ {i \in \mathbf {H}} x _ {i} \leqslant \sum_ {i \in \mathbf {I}} x _ {i},
\]

且等号仅当 \(x_{i} = 0\) 对一切 \(i \in \mathbb{C}\mathbb{H}\) 成立时才成立。

